\documentclass{article}

% Official NeurIPS style. This is the 2023 file (the only one reachable from here); the
% NeurIPS 2026 template uses `neurips_2026.sty` with an identical interface, so switching is
% a one-line change once that file is to hand.
%
% No option is passed on purpose: `final` and `preprint` must both be omitted at submission
% time, and omitting them is what produces the anonymised layout.
\usepackage{neurips_2023}

\usepackage[utf8]{inputenc}
\usepackage[T1]{fontenc}
\usepackage{hyperref}
\usepackage{url}
\usepackage{booktabs}
\usepackage{amsfonts}
\usepackage{amsmath}
\usepackage{amssymb}
\usepackage{nicefrac}
\usepackage{microtype}
\usepackage{graphicx}

\title{Learning an exact diffusion denoiser with thirteen parameters}

\author{%
  Anonymous Author(s) \\
  Affiliation \\
  Address \\
  \texttt{email}
}

\begin{document}

\maketitle

\begin{abstract}
The exact score of a diffusion model is unavailable for essentially every data distribution
of interest, which makes it impossible to separate a network's approximation error from the
intrinsic difficulty of the score it is approximating. We study a data model where the exact
score \emph{is} computable: a first-order Markov chain with continuous, non-Gaussian
innovations. Coordinatewise noising leaves the posterior a chain, chains are trees, and
sum-product belief propagation on a tree is exact --- so BP returns the Bayes-optimal
denoiser at every noise level. We then remove the assumption that the chain is known and
estimate its transition kernel by expectation maximisation through exact BP. Three facts
make this efficient: the E-step is exact rather than variational, it compresses into a single
matrix of expected transition statistics, and Fisher's identity supplies the exact gradient
of the marginal likelihood from one BP pass, so nothing is differentiated through the
recursion. On a Laplace chain, a $13$-parameter estimator fitted to $32$ noisy sequences
attains a lower score error than a $25{,}248$-parameter denoising network fitted to $4096$
--- a factor of more than $128$ in data. We report the comparison against a
locality-respecting architecture as well, where the margin is $4.1\times$ and \emph{grows}
with data rather than closing, and we state the cost: BP inference is two orders of magnitude
slower per evaluation.
\end{abstract}

\section{Introduction}

Diffusion models are trained by denoising regression. With $x = \alpha_t a + \sqrt{\Delta_t}
z$ and $z \sim \mathcal{N}(0, I)$, the objective $\mathbb{E}\|f_\theta(x,t) - a\|^2$ is
minimised by the posterior mean $\mathbb{E}[a \mid x, t]$, from which the score follows by
Tweedie's identity,
\begin{equation}
\label{eq:tweedie}
s(x,t) \;=\; -\,\frac{x - \alpha_t\,\mathbb{E}[a \mid x, t]}{\Delta_t},
\qquad \alpha_t = e^{-t},\quad \Delta_t = 1 - e^{-2t}.
\end{equation}

There is a well-known tension in what this objective can be attaining. Its global minimiser
over all measurable functions is the \emph{empirical} score --- the score of a Gaussian
mixture centred on the training points --- and a model that attained it would reproduce its
training set and nothing else. That trained diffusion models instead generalise is a central
open question, and progress on it requires data distributions for which the \emph{exact}
score is known, so that ``the network learned the wrong score'' can be distinguished from
``this score is hard''. The available such families are essentially Gaussians and mixtures
of Gaussians, which are too structureless to be informative about how a network exploits
structure.

We use a first-order Markov chain with continuous, non-Gaussian innovations. It is a minimal
model of the temporal coherence that motivates diffusion for video and other sequential
data, and it has a property that Gaussian families do not: the score is genuinely nonlinear
in $x$, while remaining exactly computable. The reason is structural rather than analytic.
Coordinatewise noising multiplies the prior by a product of sitewise likelihood factors,
which reweights node potentials without adding edges, so the posterior is again a chain; a
chain is a tree; and sum-product belief propagation is exact on a tree. BP therefore returns
the Bayes-optimal denoiser --- not an approximation to it --- at every noise level.

Our contributions are:

\begin{itemize}
\item \textbf{An exact denoiser for a non-Gaussian, correlated data model} (\S\ref{sec:bp}),
      together with the observation that the Gaussian message closure is not merely
      approximate but is \emph{exactly} the LMMSE denoiser for any innovation law with the
      same first two moments --- so single-Gaussian ``message error'' and ``model error'' are
      the same quantity.
\item \textbf{Learning the kernel by EM through exact BP} (\S\ref{sec:em}), with three
      structural properties that make it cheap: an exact E-step, a sufficient statistic that
      compresses the whole E-step into one matrix, and a gradient from Fisher's identity that
      requires no automatic differentiation through the BP recursion.
\item \textbf{A sample-efficiency comparison against trained denoisers} (\S\ref{sec:results}),
      with the capacity, yardstick and baseline-competence controls that make it a
      measurement rather than a favourable setup --- and with the two facts that cut against
      it stated in the same place.
\end{itemize}

\section{The exact denoiser}
\label{sec:bp}

Let $a \in \mathbb{R}^n$ be drawn from $p_0(a) = \mu(a_1)\prod_{i=2}^{n} K(a_i \mid a_{i-1})$
with $K(a' \mid a) = \varphi(a' - \rho a)$ for an innovation density $\varphi$ of mean zero
and variance $q$. Noising is coordinatewise Ornstein--Uhlenbeck, so
\begin{equation}
p_t(a \mid x) \;\propto\; \mu(a_1)\prod_{i=2}^{n} K(a_i \mid a_{i-1})
\prod_{i=1}^{n} \ell_t(x_i \mid a_i),
\qquad
\ell_t(x_i \mid a_i) \propto \exp\!\Big(-\tfrac{(x_i - \alpha_t a_i)^2}{2\Delta_t}\Big).
\end{equation}
The likelihood factors are unary. They change the node potentials and leave the edge set
untouched, so the posterior factor graph is the same chain as the prior's, and the two
sum-product sweeps
\begin{equation}
L_i(a_i) \propto \!\int\! K(a_i \mid a_{i-1})\,\ell_t(x_{i-1}\mid a_{i-1})\,L_{i-1}(a_{i-1})
\,da_{i-1},
\qquad
R_i(a_i) = \!\int\! K(a_{i+1}\mid a_i)\,\ell_t(x_{i+1}\mid a_{i+1})\,R_{i+1}(a_{i+1})\,da_{i+1}
\end{equation}
give exact marginals, hence the exact posterior mean and, through \eqref{eq:tweedie}, the
exact score. We represent messages on a fixed quadrature grid; the discretisation is
spectrally accurate and its error ($9.2\times10^{-15}$ at our working resolution) is four
orders of magnitude below any learning error reported here.

\paragraph{The Gaussian closure is the LMMSE denoiser, exactly.}
A natural approximation propagates only the first two moments of each message. Because a
Gaussian message times a Gaussian likelihood is Gaussian, and the transition step maps
$\mathcal{N}(m,v)$ to a law with first two moments $(\rho m,\ \rho^2 v + q)$
\emph{regardless of the shape of $\varphi$}, moment-matched BP discards every property of the
innovation beyond its variance. It therefore returns precisely the linear minimum-mean-square
estimator $\alpha_t \Sigma_0 \Sigma_t^{-1} x$ of the covariance-matched Gaussian model. We
verify this to $6\times10^{-16}$ across five innovation families. The consequence is that at
the single-Gaussian level there is no distinction between message-representation error and
model error --- a distinction that becomes real only for richer message families, and which
we separate elsewhere.

The size of that gap locates where non-Gaussianity matters. Against exact BP, the closure's
relative score error is largest at \emph{small} $t$ and decays monotonically: for a Laplace
chain at $\rho=0.85$ it runs from $0.20$ at $t=0.08$ to $9.2\times10^{-5}$ at $t=2.4$. Large
$t$ Gaussianises everything, so the prior's shape survives only where the likelihood is weak.

\section{Learning the kernel}
\label{sec:em}

We now suppose $K$ unknown and observe only noised sequences. Writing $\theta$ for the kernel
parameters, EM alternates an E-step that computes posterior expectations under the current
$\theta$ with an M-step that maximises the expected complete-data log-likelihood. Three
properties of this instance matter.

\paragraph{The E-step is exact.} In a generic latent-variable model the E-step is variational
and the monotone-ascent guarantee is lost with it. Here the posterior is a chain for the
reason given in \S\ref{sec:bp}, so BP returns exact single-site \emph{and pairwise}
marginals. Nothing is approximated but the message representation. Empirically the
monotonicity violation is exactly zero in every run we report, which is the sharpest
available test of the whole pipeline: an error anywhere in the recursion, the pairwise
accumulation, or the M-step generically breaks it.

\paragraph{The whole E-step is one matrix.} All the M-step needs is
\begin{equation}
\Xi_{kl} \;=\; \sum_{\text{sequences}} \sum_{i=2}^{n}
\Pr\big(a_{i-1} = u_k,\ a_i = u_l \mid x\big),
\end{equation}
the continuum analogue of Baum--Welch's expected transition counts. It is sufficient and
independent of how $K_\theta$ is parameterised, so the M-step never revisits the data, its
cost is independent of the number of sequences, observations at different noise levels simply
add into one $\Xi$, and changing kernel family changes only the M-step.

\paragraph{No automatic differentiation is required.} Fisher's identity gives
\begin{equation}
\nabla_\theta \mathcal{L}(\theta) \;=\; \big\langle \Xi(\theta),\ \nabla_\theta \log K_\theta
\big\rangle,
\end{equation}
the exact gradient of the \emph{marginal} log-likelihood from a single BP pass. BP is not a
computation graph to backpropagate through; it is the oracle that supplies the expectation.
We verify this against finite differences of the exact evidence to $\sim\!10^{-9}$. The
implementation is plain numerical linear algebra throughout.

\paragraph{What noising costs.} Identifiability is never destroyed by noising: the
characteristic function satisfies $\varphi_{t,\theta}(\xi) = \varphi_\theta(\alpha_t
\xi)\exp(-\Delta_t\|\xi\|^2/2)$, the Gaussian factor has no zeros, and $\alpha_t > 0$ for
every finite $t$, so $p_{t,\theta} = p_{t,\theta'}$ forces $K_\theta = K_{\theta'}$. What
degrades is \emph{information}: the map is injective but ill-conditioned, this being Gaussian
deconvolution. Measured per sequence, the Fisher information for the innovation scale falls
by $142\times$ between $t=0.05$ and $t=1.6$, while that for the correlation falls by
$26\times$. Non-Gaussianity is the first thing the channel destroys.

\paragraph{Estimator.} We report a mixture-innovation kernel: $K(a'\mid a) = \sum_{c}
\pi_c\,\mathcal{N}(a' - \rho a;\ \mu_c, \sigma_c^2)$, which retains the linear
autoregression but leaves the innovation law essentially free. It has never seen the density
it must recover. Fitted to noisy observations only, it returns the innovation variance
essentially exactly and recovers the excess kurtosis of a Laplace innovation ($3.0$)
monotonically in the component count, reaching $2.71$ at eight components.

\section{Results}
\label{sec:results}

All comparisons use a Laplace chain, $n=32$, $\rho=0.85$, with denoisers scored against exact
BP under the true prior on held-out data. The reference itself is validated against
importance sampling to $1.1\times10^{-3}$ relative over six million samples, so the yardstick
is not assumed.

\subsection{Sample efficiency}

Table~\ref{tab:main} reports relative score error averaged over the noise schedule, over six
independent training seeds, with the test set and its reference held fixed across seeds.

\begin{table}[t]
\caption{EM through exact BP ($13$ parameters) against a denoising network ($25{,}248$),
at matched training-set size. Mean over six seeds, $\pm 1$ standard error.}
\label{tab:main}
\centering
\begin{tabular}{rccc}
\toprule
sequences & network & EM--BP & ratio \\
\midrule
32   & $0.6363 \pm 0.0043$ & $\mathbf{0.0651 \pm 0.0031}$ & $9.8\times$ \\
128  & $0.5117 \pm 0.0025$ & $\mathbf{0.0333 \pm 0.0035}$ & $15.4\times$ \\
512  & $0.2813 \pm 0.0013$ & $\mathbf{0.0211 \pm 0.0022}$ & $13.3\times$ \\
1024 & $0.2154 \pm 0.0009$ & $\mathbf{0.0182 \pm 0.0021}$ & $11.8\times$ \\
2048 & $0.1756 \pm 0.0009$ & $\mathbf{0.0142 \pm 0.0009}$ & $12.4\times$ \\
4096 & $0.1643 \pm 0.0009$ & $\mathbf{0.0124 \pm 0.0005}$ & $13.2\times$ \\
\bottomrule
\end{tabular}
\end{table}

EM--BP on $32$ sequences ($0.0651 \pm 0.0031$) is better than the network on $4096$
($0.1643 \pm 0.0009$) by a factor of $2.5$ with non-overlapping intervals, so the network
requires more than $128\times$ the data to match it --- a lower bound, since $32$ was the
smallest budget tried and the network never caught up.

\paragraph{Three controls.} \emph{Capacity is not the explanation}: over $24$ configurations
from $3.2$k to $297$k parameters the best network reaches $0.208$ against EM--BP's $0.022$,
and error \emph{degrades} for the largest models, so the gap is overfitting at this data
budget rather than a badly chosen size. \emph{The baseline can learn}: trained on a Gaussian
chain where the exact denoiser is closed-form and linear, given four times the data and twice
the steps, it reaches $0.099$ --- the implementation is sound and the gap is a
sample-efficiency effect. \emph{Both parameterisations are trained}: predicting the noise wins
at low $t$ and predicting the clean signal at high $t$, by large factors, so we report the
better of the two at every level and the result rests on neither.

\subsection{What a structured architecture recovers}

The comparison above is against a fully connected network carrying no locality prior, which
is the natural first baseline but not the strongest one. We repeat it against a
weight-shared window predictor --- a one-dimensional convolution whose receptive-field radius
and parameterisation are both chosen by oracle --- over four seeds.

\begin{table}[t]
\caption{Against a locality-respecting architecture. Mean relative score error over five
noise levels, four seeds, $\pm 1$ standard error.}
\label{tab:arch}
\centering
\begin{tabular}{rccc}
\toprule
sequences & EM--BP (13) & global MLP (25{,}505) & local CNN (6{,}081) \\
\midrule
128  & $\mathbf{0.0368 \pm 0.0067}$ & $0.3595 \pm 0.0044$ & $0.0774 \pm 0.0024$ \\
512  & $\mathbf{0.0174 \pm 0.0017}$ & $0.1794 \pm 0.0035$ & $0.0570 \pm 0.0007$ \\
2048 & $\mathbf{0.0123 \pm 0.0009}$ & $0.1246 \pm 0.0027$ & $0.0505 \pm 0.0015$ \\
\bottomrule
\end{tabular}
\end{table}

Roughly $60\%$ of the deficit of the fully connected network was its architecture rather than
the estimator: the margin falls from ${\sim}10\times$ to $2.10 \pm 0.39$ at $128$ sequences.
But the remaining margin \emph{grows} with data --- $2.10 \to 3.27 \to 4.11 \pm 0.32$ --- and
the mechanism is visible in the columns. EM--BP falls by a factor of $3.0$ over a $16\times$
increase in data while the CNN falls by $1.5$ and is flattening. The structured network
converges towards a floor set by its approximation error; the estimator keeps paying down
parametric error. \textbf{The natural expectation that more data closes the gap is wrong
here}, and the honest headline against a well-chosen architecture is $4.1\times$ and rising.

\subsection{The advantage is tied to a small parametric dimension}

Replacing the continuous chain by a discrete alphabet removes every discretisation caveat at
once --- BP becomes exact up to round-off, verified against explicit enumeration to
$4.4\times10^{-15}$ --- and reproduces the same $\ge 64\times$ data-efficiency gap with zero
grid error. It also shows the limit of the effect. Sweeping the alphabet size, the advantage
over the best network at $512$ sequences falls from $15.4\times$ at six parameters to
$4.4\times$ at $132$. Nothing here suggests a structured estimator retains its advantage once
its parameter count approaches the network's.

\section{Limitations}

\textbf{Inference cost is the honest loss.} Grid BP is $O(nM^2)$ per sequence per evaluation
against a few matrix products for a network; we measure a $211$--$320\times$ slowdown.
Because reverse diffusion calls the denoiser at every integration step, this is the weakest
point of the result, and distillation of the fitted denoiser --- not attempted here --- is
the obvious remedy. \textbf{Training is slower too}: EM's cost grows linearly in the number
of sequences while the network's is fixed by its step count, so the advantage is in data, not
wall-clock. \textbf{Everything rests on the chain assumption}: under a non-Markov prior the
estimator is misspecified in a way no amount of data repairs. \textbf{One family, one length,
one correlation.} \textbf{And the comparisons above are pointwise}, whereas a diffusion model
is judged on what it generates; score error and sample quality are related by an integration
that accumulates error unevenly across the schedule, and we regard the sampling comparison as
necessary to complete the claim rather than as a formality.

\section{Conclusion}

A data model whose exact score is computable, and non-trivially so, makes it possible to ask
what a denoising network's error is made of. On this model an estimator with thirteen
parameters, learned without ever seeing a clean sample and without differentiating through
the inference algorithm, reaches an accuracy that a conventional denoiser does not reach with
two orders of magnitude more data --- and retains a growing advantage even against an
architecture built with the right inductive bias. The cost is inference time, and the scope
is a single structured family. We take the useful reading to be less that structure beats
scale, and more that when the structure of the data is known, the quantity that has to be
estimated can be very small indeed.

\end{document}
